% ============================================================
% ch17b.tex —— 第 17 章后半（17.2.3 起）
% 印刷 pp.227–250（PDF p267–p290，24 页），无插图
% 底本：Springer LNM 2386, Ch.17 Four-dimensional supersymmetric quantum field theories
% 续块：不写 \chapter
% 编号公式：17.24–17.39（逐页核图）
% ============================================================
，一个在 $R$ 中，一个在 $\bar R$ 中，另一个在 $\mathfrak g_R\simeq\End(R)$ 或 $\mathfrak g_{\bar R}$ 中。超势 $W$ 必须是规范不变的。
% [Lead 接缝修正 2026-11-03：本句承接 ch17a p.226 末尾「…都含一个场」，为可读性把首词补为「一个在 $R$ 中」，与原文 "one in R, one in R̄, and another…" 对齐。]
% [原文 p.227]

\begin{definition}[17.26]
一般的 $4d$ $\mathcal{N}=1$ 规范-物质拉格朗日密度由上面得到的各项相加而得到：若规范群 $G$ 分解为
\[
G=\prod_{a=1}^{k}\U(1)_a\times\prod_{j=1}^{l}G_j,
\]
其中 $G_i$ 是紧单李群，则对每个阿贝尔因子或单因子，都有一项形如式 (17.15) 或式 (17.18) 的项。阿贝尔因子带有如式 (17.20) 的 Fayet--Iliopoulos 项。手征超场 $\Phi_i$ 的贡献如式 (17.12)，若它们带规范群荷，则其动能项如式 (17.22) 那样被修改。超势 $W$ 必须是 $G$-不变的解析函数；这特别意味着只有处于 $G$ 的实表示中的手征超场才能有质量。
\end{definition}

\subsection{模空间}

考虑一个 $4d$ $\mathcal{N}=1$ 的 $G$-规范理论，具有可重整的拉格朗日密度以及手征超场的典范动能项。处于 $G$ 的表示 $R$ 中的手征超场 $\Phi$ 被看作 $(\dim R)$ 个手征超场 $\Phi^i$ 构成的向量（同样地 $\bar\Phi=(\bar\Phi_i)_{i=1}^{\dim R}$）。设 $(T^a)_{a=1}^{\dim\mathfrak g}$ 为规范群 $G$ 的生成元，并设 $A$ 为 $[|1,\dim G|]$ 的子集，使得对每个 $a\in A$，生成元 $T^a$ 对应 $G$ 的一个阿贝尔因子。把拉格朗日密度按分量展开即知，场 $F_i$ 与 $D_a$ 是辅助场（即它们不对应任何动能项）：拉格朗日量中涉及 $D$ 与 $F$ 的项只有
\[
\operatorname{tr}_{\mathfrak g}(D^2)+g\sum_{a\in A}\xi^a D_a+\bar F_i F^i+g\bar\phi_i D_a(T_R^a)^i{}_j\phi^j-\frac{\partial W}{\partial\phi^i}F^i-\frac{\partial\bar W}{\partial\bar\phi_i}\bar F_i\in\mathcal{L}.
\]
$\bar F_i$、$F^i$ 与 $D_a$ 的运动方程为
\[
\bar F_i=\frac{\partial W}{\partial\phi^i},\qquad F^i=\frac{\partial\bar W}{\partial\bar\phi_i},\qquad D^a=-g\xi^a-g\bar\phi_i(T_R^a)^i{}_j\phi^j.
\]
把 $D_a$ 与 $F_i$ 消去，得到在壳拉格朗日量，其中只依赖于标量场 $\phi^i$（及 $\bar\phi_i$）的项为
\[
\mathcal{L}\ni \bar D_\mu\bar\phi\,D^\mu\phi-V(\phi^i,\bar\phi_i),
\]
其中 $V(\phi^i,\bar\phi_i)$ 是（半正定的）\emph{标量势}：
% [原文 p.228]
\begin{equation}
V(\phi^i,\bar\phi_i)=\frac{\partial W}{\partial\phi_i}\frac{\partial\bar W}{\partial\bar\phi_i}+\frac{g^2}{2}\sum_a\Bigl|\bar\phi_i(T^a)^i{}_j\phi^j+\xi^a\Bigr|^2,
\tag{17.24}
\end{equation}
其中我们约定：只要 $a\notin A$ 就取 $\xi^a=0$。

% [Tutor Note: 式 (17.24) 与上一页运动方程中的指标位置不同：EOM 写 $\bar F_i=\partial W/\partial\phi^i$、$F^i=\partial\bar W/\partial\bar\phi_i$，而 (17.24) 写 $\partial W/\partial\phi_i$、$\partial\bar W/\partial\bar\phi_i$。核对方式：$\phi_i$ 与 $\phi^i$ 由典范度量/Killing 型互相升降，(17.24) 的第一项必须理解为内积 $\partial W\,\overline{\partial W}$；两处写法只是同一内积的等价表示，若当作两个独立量会得到不实的因子。]

\begin{definition}[17.27]
量子场论（QFT）的\emph{真空}是一个洛伦兹不变态，其中所有场的动能都为零。洛伦兹不变意味着只有洛伦兹不变的算子才能有非零的\emph{真空期望值}（VEV），记作 $\langle\phi^i\rangle$。若相应的 VEV 集合使标量势 $V$ \emph{局部}取极小，则称该真空是\emph{亚稳的}；如果它全局取极小，则称它是\emph{稳定的}。
\end{definition}

一个亚稳真空保持超对称，当且仅当所有超荷都平凡地作用于它。回顾式 (17.2)，这意味着在真空期望值（VEV）处求值的标量势 $V$ 为零。因此，任何亚稳的超对称真空都是真正稳定的，因为势 $V$ 半正定。

式 (17.24) 表明，$4d$ $\mathcal{N}=1$ 规范理论的超对称真空与满足 \emph{F 项}与 \emph{D 项}方程的标量 VEV 集合一一对应：
\[
\bar F_i(\phi)=0,\qquad D^a(\phi,\bar\phi)=0.
\]
规范等价的真空必须在物理上等同，这引出如下定义：

\begin{definition}[17.28]
（超对称真空的）\emph{经典模空间}记作 $\mathcal{M}_{\mathrm{cl}}$，它是标量真空期望值（VEV）构成的复仿射空间内，由 F 项与 D 项方程的解、模去规范对称性所得的子流形。更形式化地说，$\mathcal{M}_{\mathrm{cl}}$ 可描述为 F 项方程的解集按群 $G$ 所作的凯勒商。另一种方式是，用几何不变量理论定义 $\mathcal{M}_{\mathrm{cl}}$：它是由规范不变标量算子的 VEV 参数化的空间，并计入经典关系与 F 项。这两种进路的等价性由 Kempf--Ness 定理给出；另见 [252]。
\end{definition}

\subsection{四维 $\mathcal{N}=1$ SQCD（I）}

我们以 $\mathcal{N}=1$ SQCD（超量子色动力学，Super Quantum ChromoDynamics）这一例子来结束本节，它是量子色动力学的 $\mathcal{N}=1$ 推广，也是一般讨论的丰富示例。这里我们把 SQCD 当作经典 $4d$ $\mathcal{N}=1$ 场论来考虑，把量子效应的研究留到下一节。

SQCD 是一个（可重整的）拉格朗日规范理论，规范群为 $\SU(N)$（$N$ 是色数），含有 $F>1$ 对手征超场 $(Q_i,\tilde Q^i)$（$i=1,\dots,F$），称为\emph{味}，使得每个 $Q_i$（相应地 $\tilde Q^i$）在 $\SU(N)$ 的基本（相应地反基本）表示中变换，且没有超势。显式地，在 $\mathcal{N}=1$ 超空间中，拉格朗日量写作
% [原文 p.229]
\[
\mathcal{L}=\frac{1}{16\pi}\operatorname{Im}\Bigl(\tau\int\dd^2\theta\,\operatorname{tr}W^\alpha W_\alpha\Bigr)+\int\dd^2\theta\,\dd^2\bar\theta\,\Bigl(\bar Q_i \ee^{2gV}Q_i+\bar{\tilde Q}^i\ee^{-2gV}\tilde Q^i\Bigr).
\]
$\mathcal{L}$ 的单因子或阿贝尔（0-形式）整体对称性因子为：
\begin{equation}
\begin{array}{c||ccccc}
 & \SU(F)_L & \SU(F)_R & \U(1)_B & \U(1)_A & \U(1)_R\\ \hline
Q_i & \overline F & - & 1 & 1 & 1\\
\tilde Q^i & - & F & -1 & 1 & 1\\
\lambda & - & - & - & - & 1
\end{array}
\tag{17.25}
\end{equation}
如前，$\lambda$ 是\emph{规范微子}（gaugino），即矢量超多重态中的实旋量。在 $R$ 对称性下 $\lambda$ 的荷必须为 $1$（这来自 $W_\alpha\sim-i\lambda_\alpha+\dots$ 以及 $R[\dd\theta]=-1$），而 $Q_i$ 与 $\tilde Q^i$ 的 $R$ 荷必须相等，因为 $Q_i$ 与 $\tilde Q^i$ 互为电荷共轭（电荷共轭与超荷可交换）。注意把 $\U(1)_A$ 与 $\U(1)_R$ 组合起来会得到其他（虽不独立）的 $R$ 对称性，在它们之下 $\lambda$ 的荷恒为 $1$，而 $Q_i$ 与 $\tilde Q^i$ 可以带任意荷 $a\in\R$。群 $\SU(F)_L\times\SU(F)_R$ 是 SQCD 的\emph{味对称群}。味指标记作 $i,j,\cdots=1,\dots,F$，色指标记作 $a,b,\cdots=1,\dots,N$。

设 $q^a_i(x)$（相应地 $\tilde q^i_a(x)$）是手征超场 $Q^a_i(x,\theta,\bar\theta)$（相应地 $\tilde Q^i_a(x,\theta,\bar\theta)$）的标量分量。$q^a_i(x)$ 构成标量场的 $N\times F$ 矩阵 $q(x)$，而 $\tilde q^i_a(x)$ 构成 $F\times N$ 矩阵 $\bar q(x)$。在规范变换 $G(x)\in\SU(N)$ 与味对称性 $(L,R)\in\SU(F)_L\times\SU(F)_R$ 下，这些矩阵按如下方式变换：
\[
\bar q(x)\to L\,\bar q(x)\,G^\dagger(x),\qquad q(x)\to G(x)\,q(x)\,R.
\]
SQCD 的模空间由 $q$ 与 $\bar q$ 的 VEV 参数化，它们是系数为常数的 $N\times F$ 与 $F\times N$ 矩阵。这些 VEV 仍记作 $q^a_i$ 与 $q^i_a$。作用在它们之上的规范与整体对称群是 $\SU(N)\times\SU(F)_L\times\SU(F)_R$，其中 $\SU(N)$ 由所有在时空中为常数的规范变换组成。由于 SQCD 中超势为零，真空的经典模空间是 $\C^{2NF}$ 中只由 D 项定义的子簇：
\[
(D^A)^i{}_j=q^{i\dagger}_a(T^A)^a{}_b q^b_j-\tilde q^i_a(T^A)^a{}_b\tilde q^{b\dagger}_j=0,
\qquad A=1,\dots,N^2-1,
\]
并模去规范变换。如前，$T^A$ 是 $\mathfrak{su}(N)$ 在基本表示中的生成元（在反基本表示中为 $-T^A$）。等价地，经典模空间是 $\C^{2NF}$ 的子簇，其函数环由标量规范不变算子的 VEV 组成。

需要区分 $F<N$ 与 $F\ge N$ 两种情形 [4, 213, 22, 50]。

% [原文 p.230]
\begin{itemize}
\item \textbf{$F<N$。}在 $\SU(N)\times\SU(F)_R\times\U(1)_B$ 下可以假定 $q$ 分解为奇异值形式，即 $q=\operatorname{Diag}(v_1,\dots,v_F)$，对角元之外为零，且对一切 $i$ 有 $v_i\in\R_{\ge0}$。借助 $\SU(F)_L$ 还可以对 $\bar q$ 作 QR 分解，即假定 $\bar q$ 是上三角的。此时 D 项为零这一条件事实上迫使 $\bar q=q^T$。在 $q$ 的一般取值下，所有 $v_i$ 都严格为正且互不相同。这样的 VEV 选取把规范群破缺为 $\SU(N-F)$，而由（超）希格斯机制，每个被破缺的规范生成元给一个原先无质量的复标量以非零质量。因此剩余无质量复标量的个数，即经典模空间的维数，为
\[
2NF-\bigl[(N^2-1)-((N-F)^2-1)\bigr]=F^2.
\]
在此情形下参数化模空间的规范不变标量算子，是\emph{介子矩阵} $M=(m^i_j)$ 的元素：
\[
m^i_j=\bar q^i_a q^a_j,\qquad a=1,\dots,N,\quad i,j=1,\dots,F.
\]
一般地 $m$ 的秩为 $F$，这与模空间具有复维数 $F^2$ 相一致。
\item \textbf{$F\ge N$。}如前，在 $\SU(N)\times\SU(F)_R\times\U(1)_B$ 下可以假定 $q$ 是对角的，对角元为实正数 $v_1,\dots,v_N$；而在 $\SU(F)_L$ 下可以假定 $\bar q$ 是上三角的。D 项为零意味着 $\bar q$ 是对角的，其对角元为 $\bar v_1,\dots,\bar v_N$，且对一切 $i$ 有 $|\bar v_i|^2-v_i^2=\rho$，其中 $\rho$ 与 $i$ 无关。一般地此时规范群被自发地完全破缺，即模空间的复维数为
\[
2NF-(N^2-1).
\]
参数化模空间的规范不变算子是介子矩阵 $\bar q q$ 的元素，以及\emph{重子} $b$ 与\emph{反重子} $\bar b$：
\[
\begin{aligned}
b_{i_1i_2\dots i_N}&=\epsilon_{a_1a_2\dots a_N}q^{a_1}_{i_1}q^{a_2}_{i_2}\dots q^{a_N}_{i_N},\\
\bar b^{i_1i_2\dots i_N}&=\epsilon^{a_1a_2\dots a_N}\bar q^{i_1}_{a_1}\bar q^{i_2}_{a_2}\dots\bar q^{i_N}_{a_N}.
\end{aligned}
\]
矩阵 $b$ 与 $\bar b$ 完全反对称，因此总共有
\[
2\binom{F}{N}+F^2
\]
个介子与重子，其中只有 $2NF-(N^2-1)$ 个是独立的：这些关系来自定义以及 $m$、$b$ 与 $\bar b$：一个重子与一个反重子的乘积，是由（反）重子的色指标所定义的介子矩阵 $m$ 的子式。
\end{itemize}

% [原文 p.231]
\section{量子超对称规范动力学}

在量子场论中，最重要的量子效应之一是场与耦合的重整化。这一过程意味着量子场论所编码的物理依赖于观测它时的能量标度（等价地，长度标度）。我们的关注点通常是渐近自由的规范理论：它们从很高能量（紫外，UV）处的经典行为过渡到低能（红外，IR）处非常量子化、强耦合的动力学。我们特别关心低能有效物理以及真空选取对它的影响。

低能有效物理指发生在低于理论中任何非零标度的能量下的现象。在此语境下，我们区分\emph{有能隙}真空（真空态与最低非真空态之间存在显著能隙）与\emph{无能隙}真空（允许能量任意小的无质量激发）。

超对称（SUSY）量子场论中的重整化一般比非超对称理论更易处理，因为前者满足各种超对称非重整化定理。在这些理论中，任一给定超对称真空处的低能有效理论仍保持超对称。因此，在无能隙的 SUSY 真空中，若低能有效物理可由拉格朗日量描述，则它由无质量的手征超场与矢量超场组成。它们规范不变组合的标量分量参数化超对称真空的模空间 $\mathcal{M}$。

瞬子，即临界的欧氏规范场位形，是重要的非微扰量子效应，显著影响规范理论的半经典分析。反常对量子动力学的研究也至关重要。最后我们讨论全纯性——$4d$ $\mathcal{N}=1$ 理论的关键性质，它意味着某些有效耦合全纯地依赖于 UV 参数。这将是我们研究 SUSY 规范理论低能有效物理的主要工具。

量子效应可以剧烈改变拉格朗日量子场论的模空间 $\mathcal{M}$ 的结构。因此人们区分经典模空间 $\mathcal{M}_{\mathrm{cl}}$ 与真正的量子模空间 $\mathcal{M}$。$\mathcal{M}$ 的度量、拓扑甚至维数都可能与 $\mathcal{M}_{\mathrm{cl}}$ 不同。模空间中的奇点尤为重要，因为它们标志着理论低能有效描述的失效，通常与一般有质量的场变为无质量相关。我们将深入讨论量子 $4d$ $\mathcal{N}=1$ SQCD 的具体情形，并以 Seiberg 对偶的讨论作结。

\subsection{重整化}

任何具有紧单规范群 $G$ 的四维 Yang--Mills 理论，连同带 $G$ 荷的标量与费米子，其规范耦合的单圈重整化为 [173, 304]：
% [原文 p.232]
\begin{equation}
\frac{\dd}{\dd\log\mu}g(\mu)=-\frac{g^3(\mu)}{(4\pi)^2}\Bigl[\frac{11}{3}C(\mathrm{Adj})-\frac{2}{3}n_f I(R_f)-\frac{1}{3}n_s I(R_s)\Bigr],
\tag{17.26}
\end{equation}
其中 $g(\mu)$ 是能量标度 $\mu$ 处的规范耦合，$n_f$（相应地 $n_s$）是处于表示 $R_f$（相应地 $R_s$）中的 Weyl/Majorana 费米子（相应地标量）的个数。对 $G$ 的任何表示 $R$，\emph{Casimir 数} $C(R)$ 由
\[
T_{R,a}\otimes T_R^a=C(R)\,1_R
\]
定义。它满足
\[
C(R)\dim(R)=I(R)\dim(G),
\]
其中 $I(R)$ 是表示 $R$ 的 \emph{Dynkin 指标}，其定义见式 (17.21) 之下（在此归一化下，$C(\mathrm{Adj})=I(\mathrm{Adj})$ 等于 $G$ 的对偶 Coxeter 数，例如当 $G=\SU(N)$ 时 $C(\mathrm{Adj})=N$）。

当 $n_f$ 与 $n_s$ 足够小时，非阿贝尔规范动力学意味着渐近自由；此外在一定条件下，非阿贝尔规范相互作用是渐近自由所必需的（Coleman--Gross 定理 [93]）。把式 (17.26) 改写为
\[
\frac{\dd}{\dd\log\mu}\Bigl(\frac{8\pi^2}{g^2(\mu)}\Bigr)=b_1,
\]
该微分方程的解为
\[
\frac{8\pi^2}{g^2(\mu)}=-b_1\log\frac{\tilde\Lambda}{\mu},
\]
其中 $\tilde\Lambda$ 是与 $\mu$ 无关的积分常数，称为理论的（单圈）\emph{动力学标度}。它正是规范耦合的单圈重整化发散处的标度。

以 $G$ 为规范群的拉格朗日 $\mathcal{N}=1$ 规范理论由 $G$ 的伴随表示中的一个矢量超场与处于表示 $R_\chi$ 中的手征超场 $\chi$ 组成。因此式 (17.26) 化为
\[
\frac{\dd}{\dd\log\mu}\Bigl(\frac{8\pi^2}{g^2(\mu)}\Bigr)=3C(\mathrm{Adj})-n_\chi I(R_\chi).
\]
例如，对具有 $N$ 种色（即 $G=\SU(N)$）与 $F$ 种味的 SQCD，$b_1=3N-F$。

% [Tutor Note: 用式 (17.26) 的简化形式核对 $b_1=3N-F$：SQCD 有 $2F$ 个基本/反基本手征超场（每个 $I(R)=1/2$），故 $n_\chi I(R_\chi)=2F\cdot\tfrac12=F$，于是 $b_1=3C(\mathrm{Adj})-n_\chi I(R_\chi)=3N-F$，与正文一致。]

式 (17.15) 中引入的 Yang--Mills 理论中的 theta 项是全导数，故 $\Theta$ 在微扰意义下是精确的。因此在微扰论的单圈阶可写
\begin{equation}
\frac{\dd}{\dd\log\mu}\tau=\frac{\dd}{\dd\log\mu}\Bigl(\frac{4\pi i}{g^2(\mu)}+\frac{\Theta}{2\pi}\Bigr)=-\frac{b_1}{2\pi i},
\tag{17.27}
\end{equation}
其中 $\tau$ 是式 (17.16) 中定义的复化规范耦合。该方程由下式求解：

% [原文 p.233]
\begin{equation}
\tau(\mu)=\frac{b_1}{2\pi i}\log\frac{\Lambda}{\mu},
\tag{17.28}
\end{equation}
其中与 $\mu$ 无关的\emph{全纯动力学标度} $\Lambda$ 满足 $\tilde\Lambda=|\Lambda|$。用任一标度 $\mu$ 处的耦合 $\tau(\mu)$ 表示，有
\[
\Lambda=\mu\exp\Bigl(\frac{2\pi i\tau(\mu)}{b_1}\Bigr).
\]
若一个 $\mathcal{N}=1$ 渐近自由理论在某 UV 标度 $\mu_{\mathrm{UV}}$ 处的耦合为 $\tau_{\mathrm{UV}}$，则该式给出
\[
\tau(\mu)=\tau_{\mathrm{UV}}-\frac{b_1}{2\pi i}\log\frac{\mu}{\mu_{\mathrm{UV}}},
\]
其中 $\mu$ 是任一其他能量标度，使得从 $\mu_{\mathrm{UV}}$ 到 $\mu$ 全程理论的一圈微扰近似都成立。

\subsubsection*{SUSY 非重整化定理}

SUSY 场论最重要的特征之一是：它们在重整化下的行为比一般量子场论更好。可以证明，在 $d=4$ $\mathcal{N}=1$ 规范理论中，微扰圈图修正只能生成 D 项而不能生成 F 项。特别地，超势在微扰论中是树图精确的（然而动能项是 D 项，故一般会收到微扰量子修正）。这类 SUSY 非重整化定理可以统一地从超对称场论的另一基本性质——\emph{全纯性}（将在 §17.3.4 讨论）——推出。

\subsection{瞬子}

在典范归一化下，具有紧单规范群 $G$ 的纯 Yang--Mills 理论的拉格朗日量为
\begin{equation}
\mathcal{L}[A_\mu]=-\frac{1}{2g^2}\operatorname{tr}F_{\mu\nu}F^{\mu\nu}+\frac{\Theta}{16\pi^2}\operatorname{tr}F_{\mu\nu}(\star F)^{\mu\nu}.
\tag{17.29}
\end{equation}
每个连通紧单群 $G$ 都满足 $\pi_3(G)\simeq\Z$。这来自如下事实：每个这样的 $G$ 都同伦等价于一束维数 $d_i\ge3$ 的奇数维球面之并（或再取一个有限群的商），其中恰有一个球面的维数为 $3$。这使得纯 Yang--Mills 理论可以有无穷可数个拓扑不等价的经典真空。在量子理论中会发生隧穿；理论的量子真空（称为 $\Theta$-真空）具有精细的结构。隧穿振幅可以借助\emph{瞬子}方便地计算 [348, 117, 368, 118, 345]。

% [原文 p.234]
\begin{definition}[17.29]
\emph{瞬子}是欧氏理论中作用量有限的经典解：
\begin{equation}
S_E[A_\mu]=\int_{\mathcal{M}}\dd^4x\Bigl(\frac{1}{2g^2}\operatorname{tr}F_{\mu\nu}F^{\mu\nu}-\frac{i\Theta}{16\pi^2}\operatorname{tr}F_{\mu\nu}(\star F)^{\mu\nu}\Bigr),
\tag{17.30}
\end{equation}
其中在本节中 $\mathcal{M}$ 表示经 Wick 旋转的欧氏时空。
\end{definition}

有
\begin{equation}
\begin{aligned}
\int_{\mathcal{M}}\dd^4x\,\frac{1}{2g^2}\operatorname{tr}F^2&=\int_{\mathcal{M}}\dd^4x\Bigl(\frac{1}{4g^2}\operatorname{tr}(F\mp\star F)^2\pm\frac{1}{2g^2}\operatorname{tr}F_{\mu\nu}(\star F)^{\mu\nu}\Bigr)\\
&\ge\pm\int_{\mathcal{M}}\frac{1}{2g^2}\operatorname{tr}F_{\mu\nu}(\star F)^{\mu\nu}=\pm\frac{8\pi^2}{g^2}n,
\end{aligned}
\tag{17.31}
\end{equation}
其中
\[
n=\frac{1}{16\pi^2}\int_{\mathcal{M}}\dd^4x\,\operatorname{tr}F_{\mu\nu}(\star F)^{\mu\nu}=\frac{1}{8\pi^2}\int_{\mathcal{M}}\operatorname{tr}F\wedge F
\]
是规范场位形的瞬子数，或用更数学的话说，是与 $G$-主丛（$A_\mu$ 是其上的联络）以及 $G$ 的基本表示相联系之（实）向量丛的第一 \emph{Pontryagin 数}。式 (17.31) 中的不等式在 $F=\pm\star F$ 时取等，即当 $A_\mu$ 是 $\mathcal{M}$ 上的自对偶（相应地反自对偶）$G$-联络时取等。$\mathcal{M}$ 上的自对偶（相应地反自对偶）联络满足 $n\ge0$（相应地 $n\le0$）。于是：

% [Tutor Note: 式 (17.31) 的符号约定可用自对偶情形核对：取 $F=\star F$ 时 $\operatorname{tr}(F-\star F)^2=0$，不等式取等号给出 $\int\frac{1}{2g^2}\operatorname{tr}F^2=\frac{8\pi^2}{g^2}n$；反之取 $F=-\star F$ 时用 $(F+\star F)^2$ 分支，得 $n\le0$。即 $\mp$ 的选取与自对偶性一一对应。]

\begin{definition}[17.30]
当 $n\ge0$（相应地 $n\le0$）时，瞬子数为 $n$ 的\emph{瞬子}（相应地\emph{反瞬子}）是使上述界取等的欧氏规范场位形。
\end{definition}

瞬子数为 $n\ge0$ 的位形对路径积分的贡献为
\begin{equation}
\exp(-S_E)=\exp\Bigl(-n\Bigl(\frac{8\pi^2}{g^2}-i\Theta\Bigr)\Bigr)=\exp(2\pi in\tau)=\Bigl(\frac{\Lambda}{\mu}\Bigr)^{b_1n}.
\tag{17.32}
\end{equation}
这表明当 $g^2$ 小时（等价地当 $\operatorname{Im}\tau$ 大时）瞬子贡献是指数小的；它们是非微扰效应。

\begin{remark}[17.31]
当 $G$ 单连通时，例如 $G=\SU(N)$，瞬子数 $n$ 总是整数，只要式 (17.29) 中的迹取为 $\mathfrak g=\operatorname{Lie}(G)$ 上的 Killing 型且使 $\mathfrak g$ 的长根长度平方为 $2$（例如见 [181]）。当 $G$ 不单连通时，$n$ 可以是有理数。为具体起见，我们限于 $\operatorname{Lie}(G)=\mathfrak{su}(N)$ 的群 $G$，并设 $X$ 是任一光滑欧氏四维流形。在同构意义下，$X$ 上的 $\SU(N)$-丛完全由 $n\in\Z$ 决定。相反，$X$ 上的 $\SU(N)/\Z_p$-丛（其中 $p$ 整除 $N$）由瞬子数 $n$ 以及一个离散磁通量 $w_2\in H^2(X,\Z_N)$ 决定。这两个不变量由下式相联系：
% [原文 p.235]
\[
n=\frac{N(p-1)}{2p^2}\int_X\mathcal{P}(w_2)\quad\text{mod. }1,
\]
其中 $\mathcal{P}$ 是 \emph{Pontryagin 平方}运算。使 $n$ 非整数的规范场位形称为\emph{分数瞬子}。
\end{remark}

即使 $\mathcal{M}$ 的拓扑平凡，规范理论也描述延展对象。例如，'t Hooft 线对应于奇异的规范场位形；在研究存在这样一条线时的瞬子时，应当考虑的有效拓扑不是 $\mathcal{M}$，而是 $\mathcal{M}\setminus\{\text{line}\}$，后者具有非平凡拓扑。因此，即使在平坦时空中似乎无法区分 $\SU(N)$ 规范理论与 $\SU(N)/\Z_p$ 理论，二者之别仍由理论中\emph{线算子}的谱捕捉 [7]。事实上，还存在彼此不同的 $\SU(N)/\Z_p$ 规范理论，称为\emph{整体变体}：由于 $n$ 是有理数而非整数，$\Theta$ 角一般具有大于 $2\pi$ 的周期性。在 $\SU(N)/\Z_p$ 规范理论中把 $\Theta\to\Theta+2\pi$ 往往不是理论的对称性，而是把该理论的一个整体变体映到另一个整体变体。

超对称规范理论中的瞬子与超对称及 BPS 态有密切关系，因为（反）自对偶规范场在半数超对称生成元下不变。超对称有助于计算瞬子贡献，这些贡献在 $\mathcal{N}=1$ 超 Yang--Mills 理论的低能规范微子凝聚中起作用。我们将在 §17.3.5 回到 $\mathcal{N}=1$ SYM 中的规范微子凝聚。

\subsection{反常}

\begin{definition}[17.32]
若拉格朗日场论的一个经典对称性不是相应量子理论的对称性，则称它是\emph{反常的}。
\end{definition}

若该对称性是连续的，且经典守恒的 1-形式流为 $j$，则反常意味着经典守恒方程 $\dd(\star j)=\partial_\mu j^\mu=0$ 在量子理论中不成立，或更确切地说，与经典守恒方程相应的 Ward--Takahashi 恒等式不成立。

\begin{remark}[17.33]
在量子化一个经典规范理论时，规范对称性对于得到没有负模态的量子理论至关重要：在自洽的量子场论中，规范对称性不能反常。这约束了理论的物质内容。另一方面，经典整体对称性反常并不与理论的整体自洽性冲突；它只意味着量子理论的实际对称性比人们从经典描述所期望的更小（尽管有时反常对称性的一部分会以不可逆对称性的形式复苏 [89]）。
\end{remark}

% [原文 p.236]
\subsubsection*{连续整体手征对称性的反常}

反常著名地出现在带手征费米子的规范理论中（例如 §17.2.4 的 $\mathcal{N}=1$ SQCD）；深入讨论参见综述 [191, 57]。

\begin{definition}[17.34]
若一个对称性对左手与右手 Weyl 费米子的作用不同，即相应的流（若对称性连续）是手征的，则称该对称性是\emph{手征的}。
\end{definition}

手征对称性常常反常，因为在对称变换下，右手与左手费米子路径积分测度的变换除非理论的物质内容被适当调节，否则不会相互抵消 [141]。手征流反常通常按 [1] 与 [46] 称为 Adler--Bell--Jackiw（ABJ）反常。它们在微扰论中单圈精确，并且（在 $d$ 为偶数的 $d$ 维量子场论中）只通过计算单圈 $(d/2+1)$-边形 Feynman 图得到（这是 Adler--Bardeen 定理；例如见 [2] 的叙述）。因此，在四维中，手征反常由费米子带一个整体流与两个规范流插入的三角形图计算。

在以 $G$ 为规范群的 Yang--Mills 理论中，若手征对称群 $\tilde G$ 是连续的、相应流为 $\tilde j^a$，则对处于 $\tilde G\times G$ 的表示 $(\tilde R,R)$ 中的每个 Weyl 费米子有
\begin{equation}
\partial_\mu(\tilde j^a)^\mu=\frac{1}{8\pi^2}\operatorname{tr}_{\tilde R}(T^a)\operatorname{tr}_R\{T^b,T^c\}F_b\wedge F_c,
\tag{17.33}
\end{equation}
式 (17.33) 的右端除非 $\tilde G$ 阿贝尔否则为零，即只有阿贝尔整体对称性因子可以反常。对费米子求和得到总反常 $\mathcal{A}$，当 $\tilde G=\U(1)$ 时为
\begin{equation}
\mathcal{A}=\sum_i\tilde q_i I(R_i)\qquad\text{其中}\qquad \partial_\mu\tilde j^\mu=\frac{\mathcal{A}}{8\pi^2}F_b\wedge F_b.
\tag{17.34}
\end{equation}
求和遍历理论中所有 Weyl 费米子，$\tilde q_i$ 是第 $i$ 个费米子的 $\tilde G$-荷，$I(R_i)$ 是第 $i$ 个费米子所变换的 $G$ 表示 $R_i$ 的 Dynkin 指标。由反常造成的荷破坏量为
\begin{equation}
\Delta Q=\int\dd^4x\,\partial_\mu\tilde j^\mu=\frac{2\mathcal{A}}{8\pi^2}\int\operatorname{tr}F\wedge F,
\tag{17.35}
\end{equation}
其中迹再次在基本表示中取。这意味着在瞬子数为 $n$（定义见 §17.3.2）的瞬子背景中，该对称性被破坏 $\Delta Q=2\mathcal{A}n$。这表明与参数为 $\alpha$ 的 $\U(1)$ 变换相联系的反常，相当于 Yang--Mills theta 角的一个平移 $\Theta\to\Theta+2\alpha\mathcal{A}$。若所考虑的理论在 $\Theta\to\Theta+2\pi k$（$k$ 为最小的这种整数）下不变，则 ABJ 反常意味着经典 $\U(1)$ 整体对称性中只有子群 $\Z_{2\mathcal{A}/k}$ 是量子理论真正的整体对称性。

% [原文 p.237]
\begin{example}[17.35]
考虑 $4d$ $\mathcal{N}=1$ SQCD（参看 §17.2.4）。理论的手征整体对称性列于式 (17.25)；在所有可能的 $R$ 对称性中，我们考虑一个一般的 $\U(1)^a_R$，在其下 $Q_i$ 与 $\tilde Q^i$ 的荷为 $a$。非阿贝尔整体对称性因子从不反常。$\U(1)_B$ 因子也不反常，因为 $I(R_{\text{fund}})=I(R_{\text{antifund}})=1/2$。至于 $\U(1)_A$ 与 $\U(1)^a_R$，可得
\[
\mathcal{A}_A=F,\qquad \mathcal{A}_R=F(a-1)+N.
\]
这来自 $I(\text{fund})=I(\text{antifund})=1/2$、$I(\text{adj})=N$，以及如下事实：手征超多重态中 $R$ 荷为 $a$ 的费米子具有 $R$ 荷 $(a-1)$。$\mathcal{A}_R$ 中的 $+N$ 贡献来自规范微子，它总有 $R$ 荷 $1$（这由 $\mathcal{N}=1$ 规范-物质拉格朗日密度的形式决定）并在规范群的伴随表示中变换。因此 $\U(1)_A$ 总是反常的（由于 SQCD 中 $\Theta$ 以 $2\pi$ 为周期，只有其子群 $\Z_{2F}$ 是量子理论的（普通）对称性），而 $R$ 对称性非反常仅当 $F\ne0$ 且 $a=(F-N)/F$。于是，对 $F\ne0$，量子 SQCD 的整体 0-形式对称性为：
\begin{equation*}
\begin{array}{c||ccccc}
 & \SU(F)_L & \SU(F)_R & \U(1)_B & (\Z_{2F})_A & \U(1)_R\\ \hline
Q_i & \overline F & - & 1 & 1 & (F-N)/F\\
\tilde Q^i & - & F & -1 & 1 & (F-N)/F\\
\lambda & - & - & - & - & 1
\end{array}
\end{equation*}
\end{example}

% [Tutor Note: $\mathcal{A}_R=F(a-1)+N$ 可逐项核对：$Q_i$ 与 $\tilde Q^i$ 中超场的 $R$ 荷为 $a$，其费米子分量荷为 $a-1$（因 $R[\theta]=1$，$\theta$ 贡献 $-1$），$2F$ 个基本/反基本费米子各给 $(a-1)\cdot\tfrac12$，合计 $F(a-1)$；规范微子处于伴随表示、$R$ 荷 $1$，给 $I(\mathrm{adj})=N$。两项相加即为所求。]

式 (17.35) 中瞬子数的出现暗示手征 ABJ 反常与拓扑之间的联系。可以证明，以 $G$ 为规范群并带费米子的 Yang--Mills 理论中的轴反常，是给定规范背景下 Dirac 算子指标的两倍。更一般地，手征反常体现 Atiyah--Singer 指标定理，并探测给定 $G$-主丛的瞬子数 $n$。在量子规范理论中，每一种可能的规范场位形都对路径积分有贡献，因此反常的手征对称性从不是量子理论的对称性，即使它是运动方程的对称性。

手征连续整体对称性在弯曲流形上可能反常，此时人们谈论\emph{引力手征反常}。手征流的经典守恒方程被一个类 ABJ 项修改，其中黎曼曲率 $R$ 出现在规范曲率 $F$ 的位置（且数值因子不同）。

\subsubsection*{规范反常}

正如已强调的，经典场论的规范对称性必须非反常，相应的量子场论才自洽。要求规范对称性非反常会约束量子场论的物质内容。人们

% [原文 p.238]
可以按反常来自小规范变换还是大规范变换，区分\emph{小}与\emph{大}规范反常（文献中这类反常有时称为局部与整体反常，尽管二者都对应于规范（即局部）对称性）。

\begin{proposition}[17.36]
小规范反常出现在偶数时空维数的手征理论中，类似于手征整体对称性的 ABJ 反常。它们在 $d=4$ 中由费米子带三个规范流插入的三角形图计算。可以发现，处于 Yang--Mills 规范群 $G$ 的表示 $R$ 中的单个 Weyl 费米子以如下因子贡献反常：
\[
\partial_\mu(j^a)^\mu\propto \operatorname{tr}_R(T^a\{T^b,T^c\})F_b\wedge F_c.
\]
当 $G=\SU(N)$（$N>2$）时，处于 $G$ 的基本、反基本、对称或反对称表示中的 Weyl 费米子的相对贡献分别为 $1$、$-1$、$N+4$、$N-4$。量子自洽性要求费米子内容使总反常为零。例如，只含一个反对称 Weyl 费米子的 $\SU(N)$ Yang--Mills 理论不自洽，但在含 $N-4$ 个处于 $\SU(N)$ 反基本表示的 Weyl 费米子与一个反对称费米子的 $\SU(N)$ Yang--Mills 理论中，规范反常相互抵消。
\end{proposition}

一般地，当 $R$ 为实或伪实表示时 $\operatorname{tr}_R(T^a\{T^b,T^c\})$ 为零；这特别意味着伴随表示中的 Weyl 费米子或有质量费米子从不贡献反常。有时温和的修改就足以治愈理论的小规范反常；例如 Green--Schwarz 机制通过向拉格朗日量添加项来抵消反常。

\begin{example}[17.37]
$\mathcal{N}=1$ $\SU(N)$ 超 Yang--Mills 理论若含 $N-4$ 个反基本手征超场与一个反对称手征超场，则不是反常的（规范微子处于伴随表示，故不贡献反常）。同样，在具有 $N$ 种色的 SQCD 中，处于 $\SU(N)$ 基本表示与反基本表示的手征超场一样多，因此规范对称性不反常。
\end{example}

Wess--Zumino 下降程序把 $d$ 维作用量在规范变换下的变分与一个 $(d+2)$ 维多项式联系起来，后者称为\emph{反常多项式}，它是 $(d+2)$ 维 Dirac 算子的 Atiyah--Singer 指标密度。$d$ 维规范反常与 $d+2$ 维手征反常之间的这种对应给出\emph{反常流入}机制：若把 $d$ 维理论看作生活在 $(d+2)$ 维理论（称为体）中的 $d$ 维缺陷上，则体中的手征反常抵消缺陷上的规范反常，使总反常为零。

即使一个 $G$-规范理论没有小规范反常，它也可能受大反常之苦。这些反常来自关联函数在非连续连接于恒等元的规范变换下缺乏不变性。这样的\emph{大规范变换}（在平坦时空中）只在 $\pi_4(G)\ne0$ 时存在。事实表明

% [原文 p.239]
所有紧单李群 $G$ 都有 $\pi_4(G)=0$，但 $G=\USp(N)$（$N\ge2$）例外，此时 $\pi_4(\USp(N))=\Z_2$。当 $I(R)$ 为半整数时，处于 $G$ 的表示 $R$ 中的 Weyl 费米子的路径积分测度在大规范变换下被映为自身的负值。因此，例如带一个基本 Weyl 费米子的 $\USp(N)$-规范理论在量子力学上不自洽。这称为 \emph{Witten 整体反常} [363]。

这类反常按 [99] 推广为所谓的 \emph{Dai--Freed 反常}（最近的讨论见 [159]）。当 $d$ 维时空流形 $X$ 是 $(d+1)$ 维流形 $Y$ 的边界时，配分函数在大规范变换下的反常变换可以编码为 $Y$ 上的一个拓扑量子场论。这一观点贯穿了现代对反常的理解（例如见 [270]）。

\subsubsection*{'t Hooft 反常}

正如 §17.1.3 末尾所讨论的，对量子场论中每个普通或高阶形式的整体连续或离散对称性 $G$，都可以考虑理论在非平凡 $G$-丛背景中的情形。与连续整体对称性相联系的流的经典守恒方程可以被背景场的函数修改。具体地说，若 $G^{(p)}$ 是连续整体对称性、相应的守恒流为 $j^{(p+1)}$，则 $G^{(p)}$ 的非平凡背景是一个 $(p+1)$-形式 $A^{(p+1)}$，$G^{(p)}$ 通过规范变换作用于它。若量子场论由作用量 $S$ 定义，则背景 $A^{(p+1)}$ 中的理论对应于修改后的作用量
\[
S[A^{(p+1)}]=S-i\int A^{(p+1)}\wedge\star j^{(p+1)},
\]
如式 (17.7)。

\begin{definition}[17.38]
若在背景 $A^{(p+1)}$ 存在时，流的守恒方程 $\dd\star j^{(p+1)}$ 以非平凡且不可约的方式依赖于 $A^{(p+1)}$，则称对称性 $G^{(p)}$ 有 \emph{'t Hooft 反常}。若一个连续整体对称性的守恒方程被另一个对称性的背景场以不可约方式修改，则这两个连续整体对称性有\emph{混合 't Hooft 反常}。't Hooft 反常的概念推广到离散对称性，如 §17.1.3 简要讨论的那样。
\end{definition}

在四维量子场论中，连续整体 0-形式对称性的 't Hooft 反常可以通过费米子带三个整体流插入的三角形图计算。

\begin{remark}[17.39]
整体对称性的 't Hooft 反常不是定义 17.32 意义上的反常：具有 't Hooft 反常的整体对称性仍是量子理论的真正对称性。更确切地说，'t Hooft 反常编码的是\emph{规范化（gauging）的障碍}：若整体对称性 $G$ 有 't Hooft 反常，则规范化 $G$ 会使量子理论不自洽。若两个整体对称性因子 $G_1$ 与 $G_2$ 有混合 't Hooft
% [原文 p.240]
反常而没有自身 't Hooft 反常，则可以规范化 $G_1$ 或 $G_2$，但这样做必然破坏剩下的对称性。
\end{remark}

由于 't Hooft 反常不是真正的反常，它们既不约束理论的物质内容，也不约束理论的量子整体对称性。它们的意义更多是动力学上的：

\begin{proposition}[17.40]
（'t Hooft 反常匹配条件）'t Hooft 反常是标度不变的量。
\end{proposition}

一个著名的应用是：具有连续整体 0-形式对称性 $G$ 且 't Hooft 反常非零的理论永远不能有能隙，否则低能有效理论中就不可能匹配该反常。't Hooft 反常是给定量子场论的内在特征，并为对偶性提供重要检验。

\subsection{全纯性}

考虑一个渐近自由的拉格朗日规范理论。这样的理论由给定截断标度 $\mu$ 处出现在拉格朗日量中的场与耦合（以及拉格朗日量本身）的数据定义，在该标度处量子场论弱耦合，即规范耦合 $g\ll1$。在能量标度的微小变化（即 $\mu\to\mu'$）下，场与耦合发生重整化，以使任一标度处的物理不变。耦合的变化由 Wilson 重整化群方程编码：
\[
\mu\frac{\dd}{\dd\mu}c_i=\beta(c).
\]

\begin{definition}[17.41]
标度 $\mu'<\mu$ 处的 \emph{Wilson 有效作用量}是通过积分掉动量大于 $\mu'$ 的所有自由度所得的作用量；它的裸耦合对应于标度 $\mu'$ 处的重整化耦合（例如见 [106, Sec. 3.1]）。
\end{definition}

有一个工作假设：必须能用相同的自由度描述标度 $\mu$ 与 $\mu'$ 处的理论，这也是把 $\mu'$ 限制为接近 $\mu$ 的原因（我们把这当作“接近”的定义）。

\begin{proposition}[17.42]
对 $4d$ $\mathcal{N}=1$ 规范理论，存在一个全纯重整化方案，使得标度 $\mu'$ 处的有效规范耦合与有效超势只全纯地依赖于标度 $\mu$ 处的耦合与场 [317]。
\end{proposition}

\begin{proof}
任何超势都全纯依赖于场，特别地，若 $\mu'$ 接近 $\mu$，标度 $\mu'$ 处的有效超势全纯依赖于标度 $\mu$ 处的场。所考察的耦合可以看作虚假手征超场的 VEV，而这些超场也全纯地出现在拉格朗日量中，这就证明了命题。把原理论 $T$ 的耦合提升为超场所定义的理论 $T'$ 仍
% [原文 p.241]
具有 $\mathcal{N}=1$ 超对称，并且一般具有更大的整体对称性，它被虚假 VEV 的选取破缺为理论 $T$ 的整体对称性。
\end{proof}

$4d$ $\mathcal{N}=1$ 理论的这一基本性质称为\emph{全纯性} [317, 319, 213, 22, 50]。结合对称性考量以及参数空间某些极限下重整化的给定形式，全纯性强烈地约束 Wilson 有效作用量。

\begin{example}[17.43]
考虑 $4d$ $\mathcal{N}=1$ SYM 中的全纯规范耦合 $\tau$。全纯性意味着
\[
\frac{\dd}{\dd\log\mu}\tau=-\frac{b_1}{2\pi i}+\sum_{n=2}^{\infty}f_n(\tau^{1-n})+\text{非微扰},
\]
其中对一切 $n\ge2$，$f_n(\cdot)$ 是编码 $n$ 圈阶 $\tau$ 的微扰重整化贡献的全纯函数。由于 Yang--Mills theta 项是拓扑的，$\Theta=2\pi\operatorname{Re}(\tau)$ 在微扰论中不重整化，即
\[
\frac{\dd}{\dd\log\mu}\tau=-\frac{b_1}{2\pi i}+\sum_{n=2}^{\infty}f_n((\operatorname{Im}\tau)^{1-n})+\text{非微扰}.
\]
现在 $(\operatorname{Im}\tau)^{1-n}$ 只当 $n=1$ 时才是 $\tau$ 的全纯函数，故全纯性意味着所有 $f_n$ 恒为零，且 $\tau$ 的重整化在微扰论中单圈精确：
\[
\frac{\dd}{\dd\log\mu}\tau=-\frac{b_1}{2\pi i}+\text{非微扰}.
\]
\end{example}

\begin{example}[17.44]
单个手征超场 $\Phi$ 的 Wess--Zumino 模型由弱耦合拉格朗日量
\[
\mathcal{L}=\int\dd^2\theta\,\dd^2\bar\theta\,\Phi^\dagger\Phi+\int\dd^2\theta\,\bigl(m\Phi^2+g\Phi^3\bigr)+\text{c.c.},
\]
在某个 UV 标度 $\mu_{\mathrm{UV}}$ 处定义。把质量 $m$ 与耦合 $g$ 提升为手征超场 $M$ 与 $G$，所得理论具有整体对称性 $\U(1)\times\U(1)_R$，在其下（虚假）场的荷为：
\[
\begin{array}{c|cc}
 & \U(1) & \U(1)_R\\ \hline
\Phi & 1 & 1\\
M & -2 & 0\\
G & -3 & -1
\end{array}
\]
接近 $\mu$ 的标度 $\mu'$ 处的有效超势 $W_{\mu'}$ 必须具有 $\U(1)\times\U(1)_R$ 荷 $(0,2)$，故它必具有如下形式：
% [原文 p.242]
\[
W_{\mu'}=m\Phi^2f\Bigl(\frac{G\Phi}{M}\Bigr),
\]
其中 $f$ 是全纯函数。在 $g,m\to0$ 的极限下，上述理论变为自由理论。因此 $W_{\mu'}\sim m\Phi^2+g\Phi^3$。这意味着 $f(t)=1+t$，从而在任何 $\mu'$ 处都有 $W_{\mu'}=m\Phi^2+g\Phi^3$。
\end{example}

全纯性也约束非微扰量子效应，如下所示：

\begin{example}[17.45]
SQCD 在接近定义标度 $\mu_{\mathrm{UV}}$ 的标度 $\mu'$ 处的全纯规范耦合 $\tau$ 满足
\[
\tau(\mu')=\frac{3N-F}{2\pi i}\log\frac{\Lambda}{\mu'}+f(\Lambda,\mu'),\qquad
\Lambda=\mu_{\mathrm{UV}}\exp\Bigl(\frac{2\pi i\tau(\mu_{\mathrm{UV}})}{3N-F}\Bigr),
\]
其中 $\tau(\mu')$ 表达式中的第一项是单圈微扰贡献，$f(\cdot,\mu')$ 是全纯函数。SQCD 的 theta 角 $\Theta$ 以 $2\pi$ 为周期，即 $\tau\to\tau+1$，等价地：
\begin{equation}
\Lambda\to\exp\Bigl(\frac{2\pi i}{3N-F}\Bigr)\Lambda
\tag{17.36}
\end{equation}
是理论的对称性。这迫使 $f$ 在式 (17.36) 的平移下不变。此外 $f$ 必须没有常数项，使得经典极限（$\Lambda\to0$）回到微扰单圈结果。这些要求的唯一解形如
\begin{equation}
\tau(\Lambda,\mu')=\frac{3N-F}{2\pi i}\log\frac{\Lambda}{\mu'}+\sum_{n=1}^{\infty}a_n\Bigl(\frac{\Lambda}{\mu'}\Bigr)^{(3N-F)n},
\tag{17.37}
\end{equation}
其中 $a_n$ 是复系数。求和中的贡献是非微扰的，对应于规范瞬子，参看式 (17.32)。因此 $\tau$ 的重整化在微扰论中单圈精确；此外，至少在 SQCD 弱耦合时，式 (17.37) 是全纯耦合的精确重整化。
\end{example}

% [Tutor Note: 式 (17.36) 与 $\tau\to\tau+1$ 的对应可由 $\Lambda=\mu\exp(2\pi i\tau/b_1)$（$b_1=3N-F$）直接核对：$\tau\to\tau+1$ 使 $\Lambda\to\exp(2\pi i/(3N-F))\Lambda$，正是 (17.36)。这也说明 (17.37) 中的幂次恰为 $(3N-F)n$，从而求和中的每一项在 (17.36) 下不变。]

SYM 与 SQCD 的弱耦合描述中存在反常整体对称性（SQCD 中的轴对称性 $\U(1)_A$ 与 SYM 中的 $R$ 对称性 $\U(1)_R$），这进一步约束 $\tau$ 的重整化。ABJ 反常的形式表明，反常对称性下参数为 $\alpha$ 的变换可以由 theta 角的适当平移补偿：SYM 情形为 $\Theta\to\Theta+2N\alpha$，带 $F$ 种味的 SQCD 情形为 $\Theta\to\Theta+2F\alpha$。把 $\Lambda$ 提升为虚假手征超场后，SYM 与 SQCD 的反常对称性成为扩大后理论的真正对称性，在其下 $\Lambda$ 以及式 (17.37) 中的 $\tau(\Lambda,\mu')$ 都带荷。$\tau(\Lambda,\mu')$ 正确变换的唯一方式是，对所有 $n\ge1$ 都有 $a_n=0$。这意味着单圈

% [原文 p.243]
贡献之外的高阶修正（例如瞬子修正）没有贡献。因此，全纯耦合常数的单圈跑动在量子层面是精确的，即不再有来自瞬子效应或更高圈项的修正。

\subsubsection*{全纯耦合与物理耦合}

全纯规范耦合 $\tau$ 与式 (17.18) 中 Yang--Mills 拉格朗日量的归一化相联系。这不是 Yang--Mills 理论中使用的标准（物理）归一化——在后者中 $\operatorname{tr}F_{\mu\nu}F^{\mu\nu}$ 项改为乘以 $(2g_p^2)^{-1}$，下标 $p$ 表示“物理”耦合。两种归一化通过非阿贝尔矢量超场的重标度相联系，具体地 $V\to g_pV$。

路径积分中的测度在此重标度下不变，且 $g$ 与 $g_p$ 之间的关系不是解析的。这是意料之中的，因为已知 SYM 与 SQCD 中 $g_p$ 的重整化在所有圈阶都会收到微扰贡献。例如，在具有 $N$ 种色与 $F$ 种味的 SQCD 中，物理规范耦合 $g_p$ 的重整化由 Novikov--Shifman--Vainshtein--Morozov（NSVZ）公式给出：
\begin{equation}
\frac{\dd}{\dd\log\mu}g_p=-\frac{g_p^3}{(4\pi)^2}\frac{3N-\sum(1-\gamma_\chi)}{1-Ng_p^2/8\pi^2},
\tag{17.38}
\end{equation}
其中求和遍历手征超场 $\chi$ 的集合，$\gamma_\chi$ 是 $\chi$ 的\emph{反常维数}。取 $F=0$ 即得 $\SU(N)$ SYM 中物理规范耦合 $g_p$ 的重整化。

% [Tutor Note: NSVZ 公式的单圈核对：当 $g_p\to0$ 时 $\gamma_\chi\to0$、分母 $\to1$，(17.38) 化为 $\dd g_p/\dd\log\mu=-g_p^3/(4\pi)^2\,(3N-F)$，与式 (17.26) 在基本表示归一化下的单圈系数一致（$C(\mathrm{Adj})=N$，$I(\text{fund})=\tfrac12$）。注意物理耦合与全纯耦合的系数相同只在这一极限下成立。]

$g$ 与 $g_p$ 之间的非解析变换律把 $g$ 的精确跑动与 NSVZ 公式给出的 $g_p$ 的跑动联系起来。NSVZ 公式中出现手征超场的反常维数 $\gamma_\chi$，来自如下事实：SQCD 中 $g$ 的跑动只在标度 $\mu'<\mu$ 的有效 Wilson 作用量中对手征超场取非典范动能项时才单圈精确 [31]。

\subsubsection*{一个重要警告}

SQCD 中规范耦合单圈跑动精确这一论断依赖于如下假设：能量标度 $\mu'$ 与 $\mu$ 足够接近，使得重整化可以仅用 $\mu$ 处的场表达。当 $3F<N$ 时 SQCD 渐近自由，意味着规范耦合随能量标度降低而增大。

% [Tutor Note: 原文此句「当 $3F<N$ 时 SQCD 渐近自由」与式 (17.38) 的系数矛盾：渐近自由要求 $b_1=3N-F>0$，即 $F<3N$（$3N-F>0$），而不是 $3F<N$。按 $3F<N$ 也会把共形窗口的边界弄错；核对时以 $\beta\propto-(3N-F)$ 的符号为准。]因此 SQCD 最终在低能、动力学标度 $\Lambda$ 附近变为强耦合。在强耦合下，用弱耦合的 UV 自由度描述理论不再有意义。某些在高能下无关的算子可能在重整化群（RG）流中变为相关，从而削弱弱耦合分析的有效性。

% [原文 p.244]
如果把超对称 Yang--Mills（SYM）中重整化规范耦合 $g_p$ 的微分方程按能量标度积分，会出现一个阈值，在其之下方程无解——规范耦合的 RG 流有效地停止。这当然不意味着 SYM 在更低能量标度处不存在；它只表明该方程只在理论弱耦合时才可靠。换言之，仅凭全纯性无法解决超对称规范理论的强耦合动力学。更详细的讨论见 [22, Sect. 2.4.4]。

\subsection{四维 $\mathcal{N}=1$ SQCD（II）}

尽管对一般规范理论的非微扰量子效应很难说什么，超对称——特别是全纯性——在这方面大有帮助。我们以对具有 $N$ 种色与 $F$ 种味的 SQCD 之量子动力学的极简概述来说明这一点（更多细节见例如 [317, 318, 319, 213, 300, 329, 22, 50]）。

\subsubsection*{纯 $\SU(N)$ SYM（$F=0$）}

在低能下，纯 $\SU(N)$ SYM 被认为带有质量能隙地禁闭，正如纯 $\SU(N)$ Yang--Mills 理论一样（规范微子是伴随费米子，故不能像 QCD 中的夸克那样切断通量线）。有大量证据支持这一点，包括理论的、实验的与计算（格点）方面的证据。

假定纯 $\SU(N)$ 与 $\PSU(N)$ SYM 共享同一拉格朗日密度。然而可以证明，局部物理确有质量能隙地禁闭，低能物理应当只通过一个有效超势 $W_{\mathrm{eff}}$ 依赖于 $\tau$（或更确切地说，依赖于以 $\tau$ 为 VEV 的虚假手征超场）。把 $\tau$ 提升为虚假手征超场 $T$ 后从 SYM 得到的扩大理论，具有一个非反常的 $R$ 对称性，在其下 $\exp(2\pi iT/N)$ 的荷为 $2$。这（以及量纲分析）把 $W_{\mathrm{eff}}$ 约束为如下形式：
\[
W_{\mathrm{eff}}=\tilde c_N\mu^3\ee^{2i\pi\tau/N}=\tilde c_N\Lambda^3,
\]
其中 $\tilde c_N$ 是常数。$T$ 中的辅助场作为规范微子双线性项 $\operatorname{tr}\lambda\lambda$ 的源，因此
\[
\langle\operatorname{tr}\lambda\lambda\rangle=-32\pi^2\ee^{\frac{2ik\pi}{N}}\Lambda^3,
\]
即真空中的规范微子凝聚（若 $\tilde c_N\ne0$）。把 VEV $\tau$ 赋给 $T$ 会把 $R$ 对称性破缺为 $\Z_{2N}$，而规范微子凝聚进一步把 $\Z_{2N}$ 自发破缺为 $\Z_2$。因此：

\begin{proposition}[17.46]
$\SU(N)$ SYM 有 $N$ 个孤立的、有能隙的真空，标号为 $k=0,\dots,N-1$；在第 $k$ 个真空中凝聚相的相位为 $2k\pi/N$。
\end{proposition}

% [原文 p.245]
\begin{remark}[17.47]
事实上，纯 SYM 的拉格朗日量只依赖于规范代数，并不决定规范群的全局形式。例如，纯 $\SU(N)$ 与 $\PSU(N)$ SYM 共享同一拉格朗日密度，且若它们确实禁闭，则都显示 $N$ 个有能隙真空。然而真空中有效的低能物理依赖于规范群的全局形式 [7]。
\end{remark}

\subsubsection*{带 $0<F<N$ 种味的 SQCD}

我们在 §17.2.4 中已说明，带 $0<F<N$ 种味的 SQCD 具有复维数 $F^2$ 的经典模空间，由 $F\times F$ 介子矩阵 $M$ 的元素参数化。全纯性决定了超势可能的非微扰贡献（在弱耦合下）。把 $\tau$ 提升为虚假手征超场，得到扩大理论，其中 $\U(1)_A$ 是真正的整体对称性。有效超势 $W_{\mathrm{eff}}$ 必须在规范与整体对称群 $\SU(N)\times\SU(F)_L\times\SU(F)_R\times\U(1)_A\times\U(1)_B$ 下不带荷，且在 $\U(1)_R$ 下荷为 $2$。这一点连同量纲分析与弱耦合极限的研究，把 $W_{\mathrm{eff}}$ 约束为如下形式：
\[
W_{\mathrm{eff}}=c_{N,F}\Bigl(\frac{\Lambda^{3N-F}}{\det M}\Bigr)^{\frac{1}{N-F}},
\]
其中 $c_{N,F}=N-F$ 为常数。这一非微扰贡献称为 Affleck--Dine--Seiberg（ADS）超势 [3, 4]。

考虑（经典）模空间中一点，在该点每个介子 VEV 取值 $v\sim\det M^{1/2F}$，规范群被极大破缺为 $\SU(N-F)$，且该点远离模空间原点（即 $v$ 很大），使得破缺发生在 UV SQCD 理论的弱耦合区域。在高于 $v$ 的能量下，$\SU(N)$ 规范群的全纯规范耦合 $\tau$ 满足
\[
\tau(\mu)=\frac{3N-F}{2\pi i}\log\frac{\Lambda}{\mu},
\]
其中 $\Lambda$ 是 UV 理论的动力学标度。相反，低于 $v$ 时 $\SU(N-F)$ 规范群的全纯耦合 $\tau'$ 跑动为
\[
\tau'(\mu)=\frac{3(N-F)}{2\pi i}\log\frac{\Lambda'}{\mu},
\]
其中 $\Lambda'$ 是有效 $\SU(N-F)$ SYM 理论的动力学标度。在 $\mu=v$ 处必须有 $\tau(v)=\tau'(v)$（这是标度匹配），由此推出
\[
(\Lambda')^3=\Bigl(\frac{\Lambda^{3N-F}}{\det M}\Bigr)^{\frac{1}{N-F}}.
\]
假定 SYM 带有质量能隙地凝聚，这意味着 $c_{N,F}=\tilde c_{N-F}$。

% [原文 p.246]
若把 UV 理论的某个手征超场赋予大质量 $m$，则在标度 $\mu>m$ 时理论是具有 $N$ 种色与 $F$ 种味的 UV SQCD 理论，而当 $\mu<m$ 时该有质量超场脱耦，人们实际上剩下具有 $N$ 种色与 $F-1$ 种味的 SQCD。另一个标度匹配论证表明
\[
(c_{N,F-1})^{N-F+1}=(N-F+1)\Bigl(\frac{c_{N,F}}{N-F}\Bigr)^{N-F},
\]
从而 $c_{N,F}=(N-F)c^{1/(N-F)}$，其中 $c_{N,0}=\tilde c_N$。当 $F=N-1$ 时，ADS 超势具有瞬子贡献的形式；在弱耦合下可以显式计算这样的瞬子贡献，得到 $c=1$，即 $c_{N,F}=N-F$。于是：

\begin{proposition}[17.48]
ADS 超势确实被生成：
\[
W_{\mathrm{eff}}=(N-F)\Bigl(\frac{\Lambda^{3N-F}}{\det M}\Bigr)^{\frac{1}{N-F}}.
\]
此外有 $\tilde c_N=N$，因此纯 $\SU(N)$ SYM 中的规范微子确实凝聚，如上所预期。然而必须强调，这远不是纯 SYM 禁闭的证明。
\end{proposition}

\begin{remark}[17.49]
ADS 非微扰贡献的存在对模空间有戏剧性后果：标量势（包括 $W_{\mathrm{ADS}}$ 的贡献）是规范不变标量 VEV 的严格正函数。因此（尽管存在复维数 $F^2$ 的经典模空间）量子模空间是空的：在场空间中不存在有限距离处的超对称真空。势在无穷远处趋于零；这种行为称为\emph{跑向无穷}（runaway）。
\end{remark}

\subsubsection*{带 $F=N$ 种味的 SQCD}

在这种情形以及 $F>N$ 时，不存在与所有虚假对称性及其他物理约束相容的有效超势。当 $F=N$ 时，经典模空间由介子矩阵 $M$ 与两个重子 $B$、$\tilde B$ 的 VEV 参数化，它们满足
\[
\det M=B\tilde B.
\]
量子模空间是经典模空间的形变；它仍由（其标量分量的 VEV）$B$、$\tilde B$ 以及 $M$ 的矩阵元参数化，但现在满足
\begin{equation}
\det M-B\tilde B=\Lambda^{2N}.
\tag{17.39}
\end{equation}
若把某个手征超场取为有质量，这获得 ADS 超势的唯一方式。式 (17.39) 是关于带 $F=N$ 种味的 SQCD 强耦合区域的极不平凡的陈述。

% [Tutor Note: 式 (17.39) 可与 ADS 超势互检：把某个味取大质量 $m$ 后，$F=N$ 的量子约束 $\det M-B\tilde B=\Lambda^{2N}$ 应回到 $F=N-1$ 的 ADS 超势。取 $M=\operatorname{diag}(M',m)$ 并在 $\det M\ne0$ 的分支上解出 $m$ 代入，即可看出幂次 $(\Lambda^{2N+1}/\det M')^{1}$ 与命题 17.48 中 $(\Lambda^{3N-F}/\det M)^{1/(N-F)}$ 在 $F=N-1$ 时一致。]

% [原文 p.247]
带 $F=N$ 种味的 SQCD 的量子模空间非空；然而 $\mathcal{M}$（与经典模空间 $\mathcal{M}_{\mathrm{cl}}$ 相反）在原点处没有奇点。在一般点处整体对称性 $\SU(F)_L\times\SU(F)_R\times\U(1)_B\times\U(1)_R$ 完全破缺，但存在整体对称性增强的子流形。在由 $M^i_j=\Lambda^2\delta^i_j$ 与 $B=\tilde B=0$ 定义的\emph{介子分支}上，整体对称性为 $\SU(F)_D\times\U(1)_B\times\U(1)_R$；而在由 $M=0$ 与 $B=-\tilde B=\Lambda^N$ 定义的\emph{重子分支}上，整体对称性为 $\SU(F)_L\times\SU(F)_R\times\U(1)_R$。在整个 $\mathcal{M}$ 上（甚至在规范理论强耦合的原点附近）都把介子与重子用作低能自由度是合理的假设，特别是因为 UV 中整体对称性的 't Hooft 反常与以介子、重子为自由度的介子分支、重子分支上的反常相匹配。

\subsubsection*{带 $F\ge N+1$ 种味的 SQCD}

对 $F\ge N+1$，可以证明作为凯勒流形有 $\mathcal{M}=\mathcal{M}_{\mathrm{cl}}$：超势没有非微扰贡献，也没有形变。这是由与前面各段结果相容所要求的。然而，奇点在经典模空间与量子模空间中的解释不同。

例如考虑 $F=N+1$ 时 $\mathcal{M}$ 原点的奇点。经典上它对应于胶子变为无质量，但在模空间的这个区域规范理论强耦合，胶子不适合作为低能自由度。在量子理论中可以论证 $\mathcal{M}$ 由介子与重子参数化；原点的奇点暗示新的无质量自由度的出现：在那里介子与重子之间的所有关系都平凡地成立。

当 $F\ge N+2$ 时，若假定介子与重子是低能自由度，则连续整体对称性的 't Hooft 反常不能在紫外（UV）与红外（IR）区域之间匹配。对 $F>3N$，超对称 QCD（SQCD）变为红外自由，意味着模空间 $\mathcal{M}$ 原点处的低能理论是自由的非阿贝尔规范理论。这一观察暗示，对所有 $F\ge N+2$，$\mathcal{M}$ 原点附近的低能有效理论可能处于这样的非阿贝尔库仑相 [318]。

\subsection{Seiberg 对偶}

在带 $F\ge N+1$ 种味的 SQCD 之经典精确模空间的原点，连续整体对称性 $\SU(F)_L\times\SU(F)_R\times\U(1)_B\times\U(1)_R$ 未被破缺。这暗示低能有效理论中有 $F$ 个无质量味。一般地，量子场论的低能极限是重整化群流的不动点（标度不变理论），它们可以是平凡的（即自由的）或非平凡的。

% [原文 p.248]
大多数时候，超对称标度不变的四维 $\mathcal{N}=1$ 量子场论是四维 $\mathcal{N}=1$ 超共形场论（SCFT）。

当 $N$ 很大且 $\epsilon=3-F/N\ll1$ 时，带 $N$ 种色与 $F$ 种味的 SQCD 之物理规范耦合 $g_p$ 的 NSVZ beta 函数为
\[
\beta(g_p)=-\frac{g_p^3}{(4\pi)^2}\frac{3N-F(1-\gamma(g_p^2))}{1-Ng_p^2/8\pi^2},
\]
其中
\[
\gamma(g_p^2)=-\frac{g_p^2}{8\pi^2}\frac{N^2-1}{N}+O(g_p^4)
\]
是手征超场的反常维数。值
\[
g_p^*=\frac{8\pi^2}{3}\frac{N}{N^2-1}\epsilon
\]
对应于重整化群流的一个（非平凡）不动点，即 $\beta(g_p^*)=0$。因此可以预期，当 $N$ 很大且 $\epsilon$ 很小时，SQCD 在低能下流向一个非平凡 SCFT [312]。

超共形 $R$ 对称性只能是 SQCD 的 $R$ 对称性。如任何四维 SCFT 一样，手征算子的标度维数等于其 $R$ 荷的 $3/2$；由此推出如下标度维数：
\[
\Delta(M)=\frac{3}{2}R(M)=\frac{3(F-N)}{F},\qquad
\Delta(B)=\Delta(\tilde B)=\frac{3N(F-N)}{2F}.
\]
$d=4$ 中标量场的幺正界为 $\Delta\ge1$；当 $3N/2\le F<3N$ 时介子满足该界，并在 $F=3N/2$ 时取到等号。因此 SQCD 可能不仅在 $\epsilon$ 很小时、而是在整个区间 $3N/2<F<3N$ 内流向这样的相互作用 SCFT。假定如此，就得到一个自洽且有说服力的图像。区间 $3N/2<F<3N$ 称为\emph{共形窗口}。

在带 $N+2\le F<3N/2$ 种味的 SQCD 中，人们预期模空间原点处的低能有效场论由 $F$ 种带某规范群荷的手征超场描述。一个自洽的选择是 $\SU(F-N)$ 理论，带 $F$ 个基本（相应地反基本）手征超场 $q$（相应地 $\bar q$）、一个中性手征超场 $M$ 以及树图超势 $W=\bar qMq$ [312]。$q$、$\bar q$ 与 $M$ 的整体荷为：
\[
\begin{array}{c|cccc}
 & \SU(F)_L & \SU(F)_R & \U(1)_B & \U(1)_R\\ \hline
q^i & F & - & N/(F-N) & N/F\\
\bar q_i & - & \overline F & -N/(F-N) & N/F\\
M^i_j & \overline F & F & - & 2(F-N)/F
\end{array}
\]

% [原文 p.249]
这个辅助理论记作 mSQCD。SQCD 与 mSQCD 中的整体对称性以及连续整体对称性的 't Hooft 反常相互匹配。SQCD 与 mSQCD 之间的关系称为 \emph{Seiberg 对偶}。

\begin{proposition}[17.50]
（Seiberg 对偶）对 $F\ge N+2$，SQCD（具有 $N$ 种色与 $F$ 种味）与 mSQCD（具有 $F-N$ 种色与 $F$ 种味）在低能下流向同一重整化不动点。注意 $(N,F)\to(N',F')=(F-N,F)$ 把区间 $N+2\le F\le3N/2$ 与 $F'\ge3N'$ 互换，反之亦然，并把 SQCD 的共形窗口 $3N/2<F<3N$ 映为 mSQCD 的共形窗口 $3N'/2<F'<3N'$。
\end{proposition}

于是：
\begin{itemize}
\item 当 $N+2\le F\le3N/2$ 时，模空间原点附近 SQCD 的低能有效物理更好地由相应 mSQCD 理论的低能物理解释，而后者是红外自由的：SQCD 处于\emph{自由磁相}；
\item 当 $3N/2<F<3N$ 时，在模空间原点附近，SQCD 流向一个相互作用共形理论，它等价地可描述为对偶理论 mSQCD 的低能不固定点；
\item 当 $F\ge3N$ 时，SQCD 是红外自由的。这是\emph{自由电相}。
\end{itemize}

\begin{remark}[17.51]
Seiberg 对偶是一种电磁对偶，因为它把一个理论的弱耦合区域映为另一个理论的强耦合区域：SQCD 与 mSQCD 的动力学标度 $\Lambda$ 与 $\Lambda_m$ 满足 $\Lambda^{3N-F}\Lambda_m^{2F-3N}=c$，其中 $c$ 是常数。
\end{remark}

% [Tutor Note: Seiberg 对偶的秩核对：$(N,F)\to(F-N,F)$ 后对偶理论的色数为 $F-N$，其渐近自由条件为 $F<3(F-N)$，即 $F>3N/2$；这正好与 SQCD 的共形窗口 $3N/2<F<3N$ 互换，从而在 $N+2\le F\le3N/2$ 时对偶理论红外自由（自由磁相）。可用这些不等式逐条验证命题 17.50 中的区间对应。]

尽管我们不会详细深入，但值得注意的是 Seiberg 对偶可以推广到许多其他类型的四维 $\mathcal{N}=1$ 规范理论。这包括带正交规范群的超对称 SQCD [312]，以及更广泛的箭图规范理论 [49]。在后一情形，Seiberg 对偶表现为箭图的变异，而箭图节点处规范群秩的变换规则遵循热带簇 $A$-变异规则（参看第 5 章）。

% [原文 p.250]
% 原文此页为空白页（本块末页），无正文。
